\documentclass{article}

% Keep the official style package unchanged.
\usepackage{neurips_2026}

\usepackage[utf8]{inputenc}
\usepackage[T1]{fontenc}
\usepackage{hyperref}
\usepackage{url}
\usepackage{booktabs}
\usepackage{amsfonts}
\usepackage{amsmath}
\usepackage{amssymb}
\usepackage{nicefrac}
\usepackage{microtype}
\usepackage{xcolor}
\usepackage{graphicx}
\usepackage{multirow}

\title{PSAgent: Title Placeholder}

\author{Anonymous Authors}

\begin{document}

\maketitle

\begin{abstract}
% Suggested logic:
% 1. setting and importance
% 2. current bottleneck
% 3. what PSAgent is
% 4. key components
% 5. strongest result with numbers
% 6. optional release statement
% Keep this as a single paragraph.
Abstract placeholder.
\end{abstract}

\section{Introduction}

% Paragraph 1:
% Introduce multi-stage agent systems and why end-to-end quality matters.
Introduction placeholder.

% Paragraph 2:
% State concrete bottlenecks:
% - end-to-end feedback
% - partial sharing
% - optimizer/interface confounds
% - benchmark limitations

% Paragraph 3:
% Present PSAgent in one sentence, then unpack the system.

% Paragraph 4:
% Add a compact contribution block with 3 points.

\section{Related Work}

\subsection{Multi-Stage and Structured Agent Systems}

\subsection{Agent Benchmarks and Evaluation Protocols}

\subsection{Interaction, Sharing, and End-to-End Learning}

% End each subsection with one sentence:
% "Unlike prior work, PSAgent ..."

\section{Problem Setup}

\subsection{Multi-Stage Task Formulation}

% Define the structured task, states, stages, actions, and terminal outcome.

\subsection{Shared and Unshared Path Semantics}

% Define the partial-share structure cleanly and early.

\subsection{Evaluation Protocol}

% Define regret, cost, exact match, and efficiency metrics.

\section{PSAgent Framework}

% Start with an overview paragraph that references Figure 1.

\begin{figure}[t]
    \centering
    % \includegraphics[width=0.998\textwidth]{figures/framework.pdf}
    \fbox{\rule{0pt}{1.8in}\rule{0.98\textwidth}{0pt}}
    \caption{Overview of PSAgent. The figure should ideally show the multi-stage pipeline, the partial-share mechanism, and one teaser result.}
    \label{fig:framework}
\end{figure}

\subsection{System Overview}

\subsection{Interaction or Sharing Mechanism}

\subsection{Learning or Decision Protocol}

% For each subsection:
% - explain why it is needed
% - describe how it works
% - state how it connects to the full system

\section{Experiments}

% Open the section by listing research questions.
% Suggested:
% Q1: Is PSAgent non-trivial and competitive?
% Q2: Do interaction mechanisms matter?
% Q3: Does partial sharing help?
% Q4: What is the efficiency tradeoff?

\subsection{Experimental Setup}

% Datasets / benchmarks / baselines / metrics.
% Keep detailed settings for appendix if needed.

\subsection{Main Results}

\begin{table}[t]
    \centering
    \small
    \caption{Main results. Include at least one efficiency-oriented metric in the main table.}
    \label{tab:main_results}
    \begin{tabular}{lcccc}
        \toprule
        Method & Success / EM & Cost & Regret & Efficiency \\
        \midrule
        Placeholder & -- & -- & -- & -- \\
        \bottomrule
    \end{tabular}
\end{table}

\subsection{Mechanism Analysis}

\subsection{Ablation Study}

\subsection{Efficiency Analysis}

% Efficiency should be a first-class result for this paper family.

\section{Conclusion}

% Keep this short and high-level.
% Suggested logic:
% 1. restate the problem
% 2. restate PSAgent
% 3. summarize main empirical takeaways
% 4. mention the clearest next step only if useful
Conclusion placeholder.

% References can be managed with BibTeX later.

\end{document}

